\documentclass[12pt, a4paper]{article} % Classe: article, report o book [8]

% --- Lingua e Codifica (Supporto Italiano) ---
\usepackage[utf8]{inputenc}  % Codifica input
\usepackage[T1]{fontenc}      % Codifica font
\usepackage[italian]{babel}   % Lingua italiana [2]

% --- Impaginazione e Margini ---
\usepackage[margin=2.5cm]{geometry} % Imposta margini (2.5cm su tutti i lati)

% --- Pacchetti Scientifici/Matematici ---
\usepackage{amsmath, amssymb, amsfonts, amsthm} % Matematica avanzata [2]
\usepackage{graphicx} % Per inserire immagini
\usepackage{hyperref} % Per collegamenti ipertestuali
\usepackage{booktabs} % Per tabelle professionali

% --- Formattazione Testo e Altro ---
\usepackage[document]{ragged2e}
\usepackage{xcolor}    % Per usare i colori 


\title{Maurizio Alvaro}
\date{07-01-2026}

\begin{document}

\maketitle

\section{Ruolo}
Maurizio Alvaro è un ingegnere di progetto con esperienza di più di 10 anni
all'interno dell'azienda e si occupa di :
\begin{itemize}
    \item Progettazione meccanica di macchine industriali.
    \item Creazione di disegni tecnici e specifiche costruttive.
    \item Collaborazione con il reparto Project Management per la definizione delle
          specifiche tecniche.
\end{itemize}

\section{Contatto con l'attuale programma}
Maurizio non utilizza ne consulta Access per le sue attività lavorative. Il
contatto con il prodotto dei commerciali avviene tramite documenti cartacei ed
eventualmente scansionati. Il contratto viene consultato dall'uffio tecnico per
:
\begin{itemize}
    \item Definire il venduto e le specifiche tecniche.
    \item Verificare le revisioni e le modifiche apportate post-vendita dai PE.
    \item Verificare il luogo di montaggio per:
          \begin{itemize}
              \item Richiedere verifiche tecniche di specifiche macchine \textit{(es. Globosilo
                        esterni in ambienti sismici)}
              \item Definire le specifiche di trasporto.
              \item Verificare se non esistono delle limitazione particolari all'utilizzo di
                    determinati componenti.
          \end{itemize}
    \item Verificare le condizioni di montaggio e installazione.
\end{itemize}

\section{Proposte di miglioramento}
\subsection{Approccio "macro-codici"}
\begin{itemize}
    \item Creazione di codici "contenitore" (macro-codici) utilizzati in fase di vendita
          (es. TVB generico).
    \item La distinta base di vendita è completa al 90\%.
    \item L'Ufficio Tecnico interviene successivamente per "riempire" i codici con i
          disegni specifici e i dettagli costruttivi.
\end{itemize}

\subsection{Approccio "standardizzazione a catalogo"}
\begin{itemize}
    \item Spinta verso una maggiore standardizzazione delle macchine (creazione di un
          vero catalogo interno).
    \item Utilizzo di codici già esistenti e definiti per macchine standard (es. TVB
          standard, stellari taglie 1000/1500/2000), riducendo la necessità di
          personalizzazioni continue.
    \item \textbf{Vantaggio:} Il codice di vendita corrisponde già a quello di realizzazione; riduce il carico di lavoro dell'Ufficio Tecnico e minimizza gli errori.
    \item \textbf{Gestione "Fuori Standard":} Proposta di separare nettamente la gestione dello standard (flusso automatico) dal "su misura" (gestito da un reparto dedicato).
\end{itemize}

\section{Transizione digitale e gestione documentale}
È emersa una discrepanza tra le esigenze logistiche e quelle tecniche riguardo al supporto (cartaceo vs digitale).

\begin{itemize}
    \item \textbf{Problematica attuale:} L'uso di copie cartacee multiple (Vendite, Ufficio Tecnico, Elettrico) porta a perdita di informazioni e note disallineate.
    \item \textbf{Esigenza ufficio tecnico:} Necessità di passare a un sistema digitale centralizzato.
    \item \textbf{Vantaggio del digitale:} Tutti visualizzano lo stesso documento aggiornato; le modifiche sono visibili in tempo reale a tutti i reparti, evitando errori dovuti a versioni obsolete.
\end{itemize}

\section{Gestione delle revisioni e contratti}
\begin{itemize} 
    \item \textbf{Integrazione modifiche:} Le revisioni tecniche post-vendita devono essere integrate nel contratto digitale o in una "copia tecnica" speculare.
    \item \textbf{Vincoli Legali:} Il contratto di vendita firmato (e sdoganato) non può essere modificato liberamente; pertanto, serve un documento parallelo per le specifiche tecniche aggiornate senza alterare l'accordo commerciale originale.
    \item \textbf{Tracciabilità:} Ogni modifica deve essere tracciata con data, autore e descrizione della variazione per garantire trasparenza e responsabilità.
    \item \textbf{Accessibilità:} Tutti i reparti coinvolti (Commerciale, Ufficio Tecnico, Logistica) devono avere accesso alle versioni aggiornate dei documenti rilevanti.
    \item \textbf{Revisione contrattuale:} Necessità di implementare un sistema che consenta di riproporre contratto tecnico al cliente divenendo quello ufficiale.
\end{itemize}

\section{Prossimi passi suggeriti}
\begin{enumerate}
    \item \textbf{Analisi Statistica:} Utilizzare i dati per identificare le macchine più vendute (Best Sellers).
    \item \textbf{Definizione Standard:} Concentrare gli sforzi di standardizzazione sui "Best Sellers" per creare codici univoci a catalogo.
    \item \textbf{Flusso Dati:} Definire se il flusso sarà "Riempimento Macro-codici" o "Selezione da Catalogo" (l'orientamento tecnico è per la seconda).
\end{enumerate}

\section{Gestione documentale attuale}
Attualmente, la gestione documentale avviene principalmente in formato cartaceo
, ma l'ufficio tecnico con le utlime modifiche dei flussi di lavoro non
consulta più quasi il contratto. L'analisi attuale è fatta dall'ufficio PE che
consulta il contratto cartaceo per le delineare specifiche tecniche.

\end{document}